\documentclass[11pt, a4paper]{article}

% ── Encoding & language ───────────────────────────────────────────────────────
\usepackage[T1]{fontenc}
\usepackage[utf8]{inputenc}
\usepackage[british]{babel}

% ── Page geometry ─────────────────────────────────────────────────────────────
\usepackage[a4paper, top=25mm, bottom=25mm, left=25mm, right=25mm, headheight=22pt]{geometry}

% ── Fonts ─────────────────────────────────────────────────────────────────────
\usepackage{mathptmx}
\usepackage{microtype}

% ── Colour ────────────────────────────────────────────────────────────────────
\usepackage[dvipsnames, table]{xcolor}
\definecolor{headblue}{HTML}{1F4E79}
\definecolor{ruleblue}{HTML}{2E75B6}
\definecolor{accentbg}{HTML}{D6E4F0}
\definecolor{codebg}{HTML}{F5F8FC}
\definecolor{codefg}{HTML}{1F497D}
\definecolor{gray1}{HTML}{595959}
\definecolor{lightgray}{HTML}{CCCCCC}
\definecolor{varcol}{HTML}{6B3FA0}
\definecolor{datacol}{HTML}{2E7D32}

% ── Section headings ──────────────────────────────────────────────────────────
\usepackage{titlesec}
\titleformat{\section}{\color{headblue}\large\bfseries}{\color{headblue}\thesection}{1em}{}[{\color{ruleblue}\titlerule[1pt]}]
\titleformat{\subsection}{\color{ruleblue}\normalsize\bfseries}{\color{ruleblue}\thesubsection}{1em}{}
\titleformat{\subsubsection}{\color{gray1}\normalsize\bfseries}{\color{gray1}\thesubsubsection}{1em}{}
\titlespacing*{\section}{0pt}{18pt}{8pt}
\titlespacing*{\subsection}{0pt}{12pt}{4pt}
\titlespacing*{\subsubsection}{0pt}{8pt}{2pt}

% ── Header & footer ───────────────────────────────────────────────────────────
\usepackage{fancyhdr}
\pagestyle{fancy}
\fancyhf{}
\renewcommand{\headrulewidth}{0.4pt}
\renewcommand{\headrule}{\hbox to\headwidth{\color{ruleblue}\leaders\hrule height \headrulewidth\hfill}}
\fancyhead[L]{\small\color{gray1}FBM logic}
\fancyhead[R]{\small\color{gray1}The logic underlying FreeFBM}
\fancyfoot[C]{\small\color{gray1}\thepage}

% ── Hyperlinks ────────────────────────────────────────────────────────────────
\usepackage[colorlinks=true, linkcolor=ruleblue, urlcolor=ruleblue,
  pdftitle={FBM logic --- The logic underlying FreeFBM}, bookmarks=true, bookmarksnumbered=true]{hyperref}

% ── Tables, lists, boxes ──────────────────────────────────────────────────────
\usepackage{booktabs}
\usepackage{array}
\usepackage{tabularx}
\usepackage{longtable}
\usepackage[framemethod=default]{mdframed}
\newmdenv[backgroundcolor=accentbg, linecolor=ruleblue, linewidth=1pt,
  innerleftmargin=10pt, innerrightmargin=10pt, innertopmargin=6pt, innerbottommargin=6pt,
  skipabove=8pt, skipbelow=8pt]{notebox}
\newmdenv[backgroundcolor=codebg, linecolor=lightgray, linewidth=0.6pt,
  innerleftmargin=10pt, innerrightmargin=10pt, innertopmargin=6pt, innerbottommargin=6pt,
  skipabove=8pt, skipbelow=8pt]{rationale}
\usepackage{parskip}
\setlength{\parindent}{0pt}
\setlength{\parskip}{6pt}
\setlength{\emergencystretch}{3em}
\AtBeginDocument{\sloppy}
\usepackage{enumitem}
\setlist[itemize]{noitemsep, topsep=4pt, leftmargin=1.5em}
\setlist[enumerate]{noitemsep, topsep=4pt, leftmargin=1.8em}
\usepackage{amsmath}
\usepackage{amssymb}
\usepackage{stmaryrd}
\usepackage{needspace}

% ── Macros ────────────────────────────────────────────────────────────────────
% The name of the logic, in one place.
\newcommand{\FBM}{FBM logic}
\ExplSyntaxOn
\NewDocumentCommand{\code}{m}{
  \tl_set:Nn \l_tmpa_tl {#1}
  \tl_replace_all:Nnn \l_tmpa_tl {/} {/\allowbreak}
  \tl_replace_all:Nnn \l_tmpa_tl {\_} {\_\allowbreak}
  \texttt{\color{codefg}\l_tmpa_tl}
}
\ExplSyntaxOff
% A variable, named after its type as the app names it.
\newcommand{\vr}[1]{\mathsf{\color{varcol}#1}}
% Reading words of a fact atom.
\newcommand{\rd}[1]{\text{\itshape #1}}
% A fact atom in reading form: \fa{\vr{Person} \rd{lives in} \vr{City}}
\newcommand{\fa}[1]{(#1)}
% A type atom: \ta{Student}{\vr{Person}}
\newcommand{\ta}[2]{\mathsf{#1}(#2)}
% Data operators, constants, aggregates.
\newcommand{\dop}[1]{\mathbin{\color{datacol}#1}}
\newcommand{\dfn}[1]{\mathsf{\color{datacol}#1}}
\newcommand{\dc}[1]{\text{\color{datacol}\ttfamily #1}}
\newcommand{\TC}{\mathsf{TC}}
\newcommand{\tcl}[4]{[\TC_{#1,#2}\,#3](#4)}
\newcommand{\tcs}[4]{[\TC^{*}_{#1,#2}\,#3](#4)}
\newcommand{\Ob}{\mathsf{O}}
\newcommand{\Fb}{\mathsf{F}}
\newcommand{\pop}{\mathcal{P}}
\newcommand{\sem}[1]{\llbracket #1 \rrbracket}
\newcommand{\val}{\mathit{val}}
\newcommand{\udf}{\bot}

% ─────────────────────────────────────────────────────────────────────────────
\begin{document}

\begin{titlepage}
  \vspace*{4cm}
  \begin{center}
    {\color{headblue}\fontsize{32}{38}\selectfont\bfseries FBM logic}\\[0.6em]
    {\color{gray1}\Large The Logic Underlying FreeFBM}\\[2.5em]
    \textcolor{ruleblue}{\rule{0.6\textwidth}{1.5pt}}\\[0.6em]
    {\color{ruleblue}\large\itshape Design Document\quad\textperiodcentered\quad Version~1}\\[0.6em]
    \textcolor{ruleblue}{\rule{0.6\textwidth}{1.5pt}}\\[3cm]
    {\color{gray1}November 2026\\[0.3em] Describes FreeFBM v1.3.4}
  \end{center}
\end{titlepage}

\tableofcontents
\clearpage

% ═════════════════════════════════════════════════════════════════════════════
\section{Introduction}
\label{sec:intro}

FreeFBM is built on one design principle: every constraint and every
derivation rule of a schema is a statement in a single sorted logic. The
graphical notation of ORM2 (bars, dots, circled symbols, ranges) and the
textual notations (formula constraints, derivation rules, join paths) are
shorthands for statements in that logic. Verbalisation translates the
statements into structured natural language, population checks evaluate
them, and the SQL, NORMA and mapping tools translate them further. This
document calls the logic \emph{\FBM}, after fact-based modelling and the
app.

\FBM\ is a first-order logic over the populations of a schema, with
\begin{itemize}
  \item variables sorted by object type (entity type, value type or
        objectified fact type), including intersections of types;
  \item atoms for fact types (in role order), type membership, object
        equality, and comparisons and predicates on \emph{data values};
  \item counting quantifiers $\exists^{\le n}$, $\exists^{\ge n}$,
        $\exists^{=n}$ and $\exists^{n..m}$, so that uniqueness and frequency
        constraints stay first-order;
  \item transitive closure $\tcl{v}{w}{\varphi}{s,t}$ and its
        reflexive variant $\TC^{*}$, for acyclic and strongly intransitive
        rings and for recursive rules;
  \item data terms: data values of variables, constants, a library of data
        operators, conditional terms and aggregates
        ($\dfn{Count}$, $\dfn{Sum}$, $\dfn{Avg}$, $\dfn{Min}$, $\dfn{Max}$);
  \item a deontic marker $\Ob(\varphi)$ for constraints that ought to hold.
\end{itemize}

\subsection{Who this document is for}

The reader wants to know what the logic underneath FreeFBM \emph{is}: its
syntax, its typing, its meaning over a population, how each ORM2 constraint
and rule becomes a formula, and why it was designed this way. Every rule in
this document was checked against the code of FreeFBM v1.3.4; where the code
is the reference, the module is named (Appendix~\ref{app:code}). Design
alternatives appear only as rationale, in shaded boxes like this one:

\begin{rationale}
\textbf{Rationale.} Boxes like this explain a choice and the alternatives
that were considered. They can be skipped without losing the definition.
\end{rationale}

\subsection{Related documents}

\begin{itemize}
  \item \code{docs/logic/agg-logic.tex}, \emph{An intermediate sorted logic
        for ORM2} (June 2026): the exploratory draft from which \FBM\ grew. It
        weighs options that have since been decided, and some of its notation
        differs (it writes data equality as $\doteq$ where \FBM\ writes
        $\approx$). It is kept as history; this document supersedes it.
  \item \code{docs/verbalisation/verbalisation.tex}, \emph{FreeFBM
        Verbalisation: The Pipeline from Formulas to Sentences and the Rule
        Language that Drives It}, version~10: how formulas of \FBM\ become
        sentences in seven languages. Its Section ``Input: Formulas and
        Entries'' starts where Section~\ref{sec:constraints} of this document
        ends.
  \item \code{docs/logic/schema-mappings.tex}: mappings between schemas as
        rules whose two sides are patterns, a fragment of \FBM\
        (Section~\ref{sec:mappings}).
\end{itemize}

\subsection{Where the logic sits}

\begin{center}
\begin{tabular}{l@{\quad$\longrightarrow$\quad}l}
\toprule
\textbf{Source} & \textbf{What uses the formula} \\
\midrule
graphical constraints, ranges, subtypes & Constraints tab (formulas, sentences) \\
derivation rules (types, fact types)    & population validation, derived populations \\
formula constraints, join paths         & constraint accounting, equivalence, Fold \\
mapping rules (a fragment)              & SQL (CHECKs, assertions), NORMA role paths \\
\bottomrule
\end{tabular}
\end{center}

All formulas of a schema are produced by one function, \code{formalise}
in \code{src/logic/constraintFormulas.js}, as a list of
\emph{entries}, each with a key, the element it belongs to, a kind
(``Mandatory role'', ``Ring (acyclic)'', \dots), and either a formula or a
note saying why there is none yet. Everything downstream reads these
entries; nothing downstream reads constraint types. A hand-written formula
constraint and a graphical constraint with the same content are therefore
the same statement, with the same sentence and the same check.

% ═════════════════════════════════════════════════════════════════════════════
\section{Vocabulary: what a schema provides}
\label{sec:vocab}

A schema supplies the non-logical symbols of \FBM.

\paragraph{Types.} The \emph{types} are the object types and the
objectified fact types:
\begin{itemize}
  \item an \emph{entity type} $E$, whose instances are entities;
  \item a \emph{value type} $V$, whose instances are data values; it has a
        datatype, and through it a \emph{sort} (below), or no datatype at all;
  \item an \emph{objectified fact type} $f$, which is also a type: its
        instances are the nesting instances of $f$'s facts.
\end{itemize}

\paragraph{Fact types.} A fact type $f$ has roles $1,\dots,n$, each played
by a type $\pi_f(i)$. A fact atom always lists its arguments in this role
order (\code{roles[]} in the model). The primary reading of $f$ has its
slots in the same order, which is why atoms can be written in reading form
(Section~\ref{sec:notation}). An objectified fact type $f$ also has one
\emph{link fact type} $f_i$ per role, binary, with roles ($f$, $\pi_f(i)$),
read ``$f$ involves $\pi_f(i)$'' unless it has a reading of its own. In the
model its id is \code{<factId>\_il\_<i>}; atoms may use it like any fact
type.

\paragraph{Subtype links.} A subtype link $S \sqsubseteq T$ between entity
types (or objectified fact types). The \emph{type hierarchy} of a type is
its connected component in the graph of subtype links; its \emph{roots} are
the types in it without a supertype.

\paragraph{Reference schemes.} An entity type is identified by its
preferred identifier: a \emph{reference mode} (a binary fact type
$E\ \rd{has}\ V$ with a value type $V$ and a preferred internal uniqueness
constraint on $V$'s role), a \emph{composite identifier} (an internal or
external preferred uniqueness constraint over several roles), or one
inherited from a supertype. The reference mode matters to the logic in one
way: it gives the entity a \emph{data value} (Section~\ref{sec:values}).

\paragraph{Sorts.} The sorts of data values are FreeFBM's \emph{value
kinds} (\code{VALUE\_KINDS} in \code{src/utils/datatypeValues.js}), shared
by all datatype profiles (Abstract, SQL:2016, XML Schema, JSON, TypeScript,
NORMA, PostgreSQL/SQL types):

\begin{center}
\small
\begin{tabularx}{\textwidth}{l l X}
\toprule
\textbf{Sort} & \textbf{Order} & \textbf{Values} \\
\midrule
\code{integer}, \code{decimal} & total & numbers; decimals rounded to 15 significant digits \\
\code{float} & partial & IEEE doubles with INF, $-$INF and NaN (NaN is comparable to nothing) \\
\code{string} & total & text, compared character by character (UTF-16 code units) \\
\code{boolean} & total & true/false; NORMA's Yes/No (\code{yesNo}) is the same sort \\
\code{date}, \code{datetime}, \code{time} & partial & instants with or without a timezone \\
\code{gYear}, \code{gYearMonth}, \code{gMonth}, \code{gMonthDay}, \code{gDay} & partial & XML Schema's partial dates \\
\code{duration} & partial & months and seconds \\
\code{yearMonthDuration}, \code{dayTimeDuration} & total & the two halves of a duration \\
\code{binary} & total & bytes (byte order) \\
\code{language}, \code{hexBinary}, \code{base64Binary}, \code{unordered} & none & equality only (URIs, names, \dots) \\
\bottomrule
\end{tabularx}
\end{center}

Each sort has a parser from text to value, a value equality and, where it
has an order, a comparison that may answer ``incomparable''. A date with a
timezone and one without are comparable only when they are more than 14
hours apart, as in XML Schema. Numbers \emph{widen} implicitly:
$\code{integer} \sqsubseteq \code{decimal} \sqsubseteq \code{float}$.

\begin{rationale}
\textbf{Values, not strings.} Equality and order in \FBM\ are over data
values, never over the strings that write them: \code{en-GB} equals
\code{en-gb} as a language tag, and \code{P2D} is less than \code{P10D} as a
duration. Sorts are defined once, below the datatype profiles, so every
profile gets the same operators and the same value semantics. A datatype
that maps to no sort leaves its values without a sort; they can be compared
with each other only (as numbers when both read as numbers, else as text),
and cannot be used in calculations.
\end{rationale}

% ═════════════════════════════════════════════════════════════════════════════
\section{Populations: what formulas are about}
\label{sec:populations}

A formula of \FBM\ is true or false of a \emph{population} of the schema,
the sample data of the Population tab or a population that a tool builds.

\begin{notebox}
A \textbf{population} $\pop$ gives
\begin{itemize}
  \item for every type $T$ a finite set $\pop(T)$ of instances, and
  \item for every fact type $f$ of arity $n$ a finite set $\pop(f)$ of
        $n$-tuples, the facts of $f$.
\end{itemize}
\end{notebox}

\paragraph{Instances.} An instance of a value type is a data value of its
sort, canonicalised so that value-equal spellings are one instance. An
instance of an entity type is identified by its reference values: the value
of its reference mode, or the list of the values (or instances) of its
composite identifier. An instance of an objectified fact type is identified
by the fact it nests. Two variables denote the same object only if they
lie in the same type hierarchy and have the same identifier: a
\code{Person} identified by 7 and a \code{Room} identified by 7 are
different objects.

\paragraph{Facts.} A fact of $f$ is a tuple of instances, one per role, each
an instance of the role's player. A population is a set: the same tuple
cannot occur twice.

\paragraph{Subtypes.} $\pop(S)$ and $\pop(T)$ are separate sets; that each
instance of a subtype is an instance of its supertype is the formula of the
subtype link (Section~\ref{sec:constraints}), not a property of populations.
The populations of derived types and fact types are not stored but computed
from their rules (Section~\ref{sec:derivation}).

\paragraph{Closed world.} Formulas are evaluated over the population as
given. A quantified variable of type $T$ ranges over $\pop(T)$, so every
formula has a finite, and in the absence of aggregates and closures
polynomial, evaluation. Nothing outside the population exists for a
formula: there is no domain of ``all possible persons''.

\begin{rationale}
\textbf{Why populations and not arbitrary structures.} The draft
(\code{agg-logic.tex}) started from an infinite domain of objects with
interpretations into it. For what FreeFBM does with formulas (check sample
data, compute derived populations, build counterexamples in accounting,
write SQL) the population \emph{is} the structure, and range restriction to
the population is what makes every formula decidable by evaluation. The
price is the one the draft found for computed values: a formula cannot
mention a value that is not in the population. Derivation rules with
computed roles solve that explicitly (Section~\ref{sec:factrules}).
\end{rationale}

% ═════════════════════════════════════════════════════════════════════════════
\section{Syntax}
\label{sec:syntax}

\subsection{Variables}

Every variable has a \emph{type}: an object type or an objectified fact
type. A variable may also range over the \emph{intersection} of several
types; this is how the subject of a derived subtype with several direct
supertypes is typed (written $\vr{Student\&Employee}$). Variables are
objects without names; names are given only when a formula is shown
(Section~\ref{sec:notation}).

\subsection{Formulas}

\begin{center}
\small
\begin{tabularx}{\textwidth}{l l X}
\toprule
\textbf{Form} & \textbf{Written} & \textbf{Meaning} \\
\midrule
truth & $\top$ & always true \\
fact atom & $f(x_1,\dots,x_n)$ & the tuple is a fact of $f$ (arguments in role order; $f$ may be a link fact type) \\
type atom & $T(x)$ & $x$ is an instance of $T$ \\
object equality & $x = y$, $x \ne y$ & the same object, or not \\
value comparison & $s \approx t$, $s \not\approx t$, $s < t$, $s \le t$, $s > t$, $s \ge t$ & on the data values of terms $s$, $t$ \\
data predicate & $\dfn{p}(t_1,\dots,t_k)$ & e.g.\ $\dfn{startsWith}$, $\dfn{matches}$ \\
connectives & $\neg\varphi$, $\varphi\wedge\psi$, $\varphi\vee\psi$, $\varphi\to\psi$, $\varphi\leftrightarrow\psi$ & classical, two-valued \\
quantifiers & $\forall x\,\varphi$, $\exists x\,\varphi$ & over the population of $x$'s type \\
counting & $\exists^{\le n}\bar x\,\varphi$, $\exists^{\ge n}\bar x\,\varphi$, $\exists^{=n}\bar x\,\varphi$, $\exists^{n..m}\bar x\,\varphi$ & the number of tuples $\bar x$ with $\varphi$ is $\le n$, $\ge n$, $=n$, in $[n,m]$ \\
closure & $\tcl{v}{w}{\varphi}{s,t}$, $\tcs{v}{w}{\varphi}{s,t}$ & a chain of one or more (zero or more) $\varphi$-steps from $s$ to $t$ \\
deontic & $\Ob(\varphi)$ & $\varphi$ ought to hold \\
\bottomrule
\end{tabularx}
\end{center}

The arguments of fact, type and equality atoms and of closures are
variables, never terms; data terms occur only in value comparisons and data
predicates. The left side of a value comparison is never a constant: a
comparison is written $x \approx \dc{'c'}$ or $x < y$, never
$\dc{'c'} < x$, because the datatype of the other side decides how the
constant is read.

Two further modal forms exist only inside verbalisation and are never
stored or evaluated: $\Diamond\varphi$ (``it is possible that'', for the
default verbalisation of a missing mandatory or uniqueness constraint) and
$\Fb(\varphi)$ (``it is forbidden that'', the negative form of a deontic
constraint). The negative verbalisation also uses $\exists^{>n}$ (``more than
$n$''); see the verbalisation specification.

\subsection{Data terms}

\begin{center}
\small
\begin{tabularx}{\textwidth}{l l X}
\toprule
\textbf{Term} & \textbf{Written} & \textbf{Meaning} \\
\midrule
variable & $x$ & the data value of $x$ (Section~\ref{sec:values}) \\
constant & $\dc{'Smith'}$, $\dc{12}$ & a value written as text, read under the sort its place needs \\
operator & $s \dop{+} t$, $\dfn{year}(t)$, \dots & a data operator (Section~\ref{sec:operators}) \\
conditional & $(\mathbf{if}\ \varphi\ \mathbf{then}\ s\ \mathbf{else}\ t)$ & $s$ when $\varphi$ holds, else $t$ \\
aggregate & $\dfn{Sum}\{t \mid v_1,\dots,v_k : \varphi\}$ & the sum of the bag of values of $t$, one per binding of $\bar v$ with $\varphi$ \\
count & $\dfn{Count}\{v_1,\dots,v_k : \varphi\}$ & the number of bindings of $\bar v$ with $\varphi$ \\
\bottomrule
\end{tabularx}
\end{center}

An aggregate may be marked \emph{distinct}
($\dfn{Sum}_{\mathit{distinct}}\{\dots\}$): its bag is then a set of values.
Aggregates and conditional terms contain formulas, so terms and formulas
nest in each other.

\subsection{Binding}

$\forall$, $\exists$ and the counting quantifiers bind their variables in
their body. An aggregate binds its variables in its condition and its term.
$\tcl{v}{w}{\varphi}{s,t}$ binds $v$ and $w$ in $\varphi$ only; $s$ and $t$
are outside the brackets, and every other free variable of $\varphi$ is a
\emph{parameter} of the closure. The square brackets delimit the scope of
$v$ and $w$, so they may reuse names from outside.

% ═════════════════════════════════════════════════════════════════════════════
\section{Data values of variables}
\label{sec:values}

Value comparisons, data predicates and data operators work on data values,
not on objects. A variable $x$ of type $T$ has a data value when
\begin{itemize}
  \item $T$ is a value type: $\val(x) = x$, of $T$'s sort; or
  \item $T$ is an entity type with a reference mode, its own or inherited
        from a supertype: $\val(x)$ is the value $x$ is identified by, of the
        reference value type's sort.
\end{itemize}
Any other variable (an entity type without a reference mode, one with a
composite identifier, an objectified fact type) has no data value, and
using it in a data term is a type error. A variable over an intersection
has the data value of the first of its types that has one.

In a rule the user writes the entity variable directly in the comparison;
the formula that the Constraints tab shows makes the reference mode
explicit:
\[
  \vr{Course} \approx \dc{'DB1'}
  \quad\text{is formalised as}\quad
  \exists \vr{CourseCode}\, \bigl(\fa{\vr{Course}\ \rd{has}\ \vr{CourseCode}} \wedge \vr{CourseCode} \approx \dc{'DB1'}\bigr).
\]
Inside an aggregate whose term reads the value of one of the aggregate's
own entity variables, the value variable is added to the aggregate's
variables; since a reference mode is one-to-one, the bag is the same:
$\dfn{Sum}\{\vr{CourseCode} \mid \vr{Course}, \vr{CourseCode} : \dots \wedge
\fa{\vr{Course}\ \rd{has}\ \vr{CourseCode}}\}$.

\begin{rationale}
\textbf{Two equalities.} $x = y$ compares objects; $x \approx y$ compares
data values. For two variables of one value type they coincide, since a
value type's instances are its values. They differ for entity types
(two \code{Course} objects with value-equal codes are the same object only
because the reference mode is unique) and across value types
($\vr{ShoeSize} \approx \vr{Age}$ may hold, $\vr{ShoeSize} = \vr{Age}$
never does). The draft wrote $\doteq$ for data equality; \FBM\ uses
$\approx$, $\not\approx$ and the four order symbols, all primitive, since
with partial orders $\neg(x < c)$ does not give $x \ge c$. The draft's
implicit coercion of an entity variable to its reference value is kept for
writing rules, but every displayed and evaluated formula spells it out.
\end{rationale}

% ═════════════════════════════════════════════════════════════════════════════
\section{Typing}
\label{sec:typing}

A formula that is stored in the model (a derivation rule, a formula
constraint, a join path written as a formula) is checked before it is
formalised or evaluated (\code{checkFormula} in
\code{src/logic/derivationRule.js}, \code{typeComparison} and
\code{typePredicate} in \code{src/logic/dataOperators.js}). The check gives
\emph{errors}, which keep the formula from being used, and
\emph{warnings}, which leave a well-defined formula that is probably not
what was meant.

\subsection{Objects in places}

A \emph{place} is a role of a fact atom, the argument of a type atom, a side
of an object equality, or an end of a closure. A variable of types $\bar T$
fits a place for instances of $P$ in one of four ways (\code{typeGraph}):
\begin{center}
\small
\begin{tabularx}{\textwidth}{l X X}
\toprule
\textbf{Fit} & \textbf{When} & \textbf{Result} \\
\midrule
ok & some $T_i$ is $P$ or a subtype of $P$ & fine \\
overlap & $\bar T$ and $P$ share a supertype & fine: whether $x$ is there depends on the population \\
disjoint & they share one, but lie below mutually exclusive subtypes & warning: always false \\
none & they lie in different type hierarchies & warning: always false \\
\bottomrule
\end{tabularx}
\end{center}
Mutually exclusive subtypes are those of an exclusion or exclusive-or
constraint over subtype links. An atom that can never hold is allowed: it
is false, and the warning says that the formula can be simplified. Only the
ends of a closure are strict: an end from another hierarchy is an error.

\begin{rationale}
\textbf{Soft typing.} Variables are sorted, but a placement that cannot
hold is not a syntax error. A strict regime (a \code{Person} variable only
in \code{Person} roles) would reject the common case of a supertype variable
in a subtype's role, which simply restricts it; rejecting only the
hopeless cases (no common supertype, exclusive subtypes) and still giving
them a meaning (false) keeps imported NORMA rules and half-edited rules
usable.
\end{rationale}

\subsection{Data terms}

Every data term gets a sort.
\begin{itemize}
  \item A variable has the sort of its data value; a variable without one is
        an error.
  \item A constant has no sort of its own: it is read under the sort its
        place needs and must parse there (\dc{'P1X'} is not a duration).
        A constant that nothing around it types is guessed (integer,
        decimal, boolean, else string).
  \item An operator chooses the first of its signatures that its arguments
        fit, with numeric widening; a constant fits a parameter it parses
        under.
  \item A conditional term has the join of its branches' sorts.
  \item $\dfn{Count}$ gives \code{integer}; $\dfn{Sum}$ keeps the numeric sort
        of its term; $\dfn{Avg}$ gives \code{decimal} (\code{float} for
        floats); $\dfn{Min}$ and $\dfn{Max}$ keep any ordered sort. An
        aggregate needs at least one variable.
\end{itemize}
A value comparison needs the two sides to have a common sort (their join
under widening), and the order symbols need a sort with an order: values of
an unordered sort can only be compared with $\approx$ and $\not\approx$.
Values without a datatype may be compared with each other, variable with
variable, but not used in calculations.

% ═════════════════════════════════════════════════════════════════════════════
\section{Semantics}
\label{sec:semantics}

Let $\pop$ be a population and $\beta$ a binding of the free variables to
objects. $\pop,\beta \models \varphi$ is defined as usual for first-order
logic, with these cases:

\begin{itemize}
  \item $\pop,\beta\models f(x_1,\dots,x_n)$ iff
        $(\beta(x_1),\dots,\beta(x_n)) \in \pop(f)$, where $\pop(f)$ is the
        stored or derived population of $f$.
  \item $\pop,\beta\models T(x)$ iff $\beta(x) \in \pop(T)$.
  \item $\pop,\beta\models x = y$ iff $\beta(x)$ and $\beta(y)$ are the same
        object (same hierarchy, same identifier).
  \item $\pop,\beta\models \forall x\,\varphi$ iff
        $\pop,\beta[x\mapsto o]\models\varphi$ for every $o\in\pop(T)$, $T$
        the type of $x$ (for an intersection, the objects all its types
        share); $\exists$ likewise.
  \item $\pop,\beta\models \exists^{\le n}\bar x\,\varphi$ iff the number of
        tuples $\bar o$ over the populations of $\bar x$'s types with
        $\pop,\beta[\bar x\mapsto\bar o]\models\varphi$ is at most $n$; and so
        on for $\ge n$, $=n$ and $n..m$ (at least $n$ and at most $m$).
  \item $\pop,\beta\models \tcl{v}{w}{\varphi}{s,t}$ iff there are
        $o_0=\beta(s), o_1, \dots, o_k=\beta(t)$ with $k\ge 1$ and
        $\pop,\beta[v\mapsto o_{i-1}, w\mapsto o_i]\models\varphi$ for each
        $i$; for $\TC^{*}$ also when $\beta(s)=\beta(t)$ ($k=0$).
  \item $\Ob(\varphi)$ is evaluated as $\varphi$. The marker changes what a
        false value means, not the value: a population that breaks an
        alethic constraint is inconsistent, one that breaks a deontic
        constraint is reported as breaking an obligation.
\end{itemize}

\subsection{Data terms and undefinedness}

A data term denotes a value of its sort, or nothing ($\udf$): data
operators are partial: division by zero, an integer overflow, a text that
does not read as a date, or a day that does not exist in
$\dfn{makeDate}$ give no value. The rules are:
\begin{enumerate}
  \item A term with an undefined argument is undefined.
  \item A value comparison or data predicate with an undefined term is
        \emph{false}.
  \item $s < t$ (and $\le$, $>$, $\ge$) is false when $s$ and $t$ are
        incomparable under a partial order (NaN, a date with and one without
        a timezone close together).
  \item $(\mathbf{if}\ \varphi\ \mathbf{then}\ s\ \mathbf{else}\ t)$ is $s$
        or $t$ as $\varphi$ is true or false.
\end{enumerate}
Formulas stay two-valued: undefinedness stops at the atom. One consequence
is that $\neg(x \approx t)$ and $x \not\approx t$ differ when $t$ is
undefined: the first is true, the second false.

\subsection{Aggregates}

For $\mathcal{F}\{t \mid \bar v : \varphi\}$ under $\beta$, let $B$ be the
bag holding $\sem{t}$ for every binding $\bar o$ of $\bar v$ over their
types' populations with $\pop,\beta[\bar v\mapsto\bar o]\models\varphi$;
with \emph{distinct}, value-equal elements are kept once.
\begin{itemize}
  \item $\dfn{Count}\{\bar v:\varphi\}$ is the number of such bindings (0 when
        there are none).
  \item \emph{Strict}: if $t$ is undefined for any binding, the aggregate is
        undefined, with distinct too.
  \item $\dfn{Sum}$ of an empty bag is 0; $\dfn{Avg}$, $\dfn{Min}$ and
        $\dfn{Max}$ of an empty bag are undefined.
  \item $\dfn{Min}$ and $\dfn{Max}$ are undefined when two elements are
        incomparable.
\end{itemize}

\begin{rationale}
\textbf{Partial operators, strict aggregates, false atoms.} The
alternatives were a three-valued logic (as SQL's) or total operators with
error values. A two-valued logic in which only atoms can fail keeps every
connective classical, which the verbalisation, the negative (counterexample)
rewrite and the chase in constraint accounting rely on. Strictness makes an
aggregate over partly undefined values undefined rather than silently
leaving them out. There is no $\dfn{today}$ or $\dfn{now}$: a formula's truth
depends on the population only.
\end{rationale}

\begin{rationale}
\textbf{Labelled nulls.} Constraint accounting and the contract check build
populations with \emph{labelled nulls} $\bot_1, \bot_2,\dots$ (unknown
values that the chase may later identify). The evaluator treats a labelled
null as a value equal only to itself: $\approx$ and $\not\approx$ compare
labels, and order comparisons with a labelled null are false. These are
values of the evaluator, not of the logic; no formula can write one.
\end{rationale}

% ═════════════════════════════════════════════════════════════════════════════
\section{Data operators}
\label{sec:operators}

The data operators are a fixed library (\code{OPS} and \code{PREDS} in
\code{src/logic/dataOperators.js}), the same for every datatype profile.
Each has one or more signatures; the first that fits is used.

\begin{center}
\small
\begin{tabularx}{\textwidth}{l X}
\toprule
\textbf{Group} & \textbf{Operators} \\
\midrule
Numbers & $\dop{+}$, $\dop{-}$, $\dop{\times}$ (same numeric sort); $\dop{\div}$ (decimal or float); $\mathbin{\mathrm{div}}$, $\mathbin{\mathrm{mod}}$ (integers); unary $-$; $\dfn{abs}$; $\dfn{least}$, $\dfn{greatest}$ (any ordered sort); $\dfn{round}(x,n)$ (keeps $x$'s sort, half away from zero); $\dfn{roundInt}$, $\dfn{floor}$, $\dfn{ceiling}$, $\dfn{truncate}$; $\dfn{power}$, $\dfn{sqrt}$ (floats) \\
Text & $\dop{\Vert}$ (concatenation), $\dfn{length}$, $\dfn{substring}(s, \mathit{start}, \mathit{length})$, $\dfn{upper}$, $\dfn{lower}$, $\dfn{trim}$, $\dfn{toText}$ \\
Conversions & $\dfn{asInteger}$, $\dfn{asDecimal}$, $\dfn{asFloat}$, $\dfn{asBoolean}$, $\dfn{asDate}$, $\dfn{asDatetime}$, $\dfn{asTime}$: a text (or a narrower value) read as another sort, partial \\
Dates and times & $\dfn{year}$, $\dfn{month}$, $\dfn{day}$, $\dfn{hour}$, $\dfn{minute}$, $\dfn{second}$, $\dfn{dayOfWeek}$ (1 = Monday), $\dfn{yearsBetween}$, $\dfn{makeDate}$, $\dfn{dateOf}$, $\dfn{timeOf}$; $+$ and $-$ of a duration; the difference of two instants as a day-time duration \\
Durations & $+$, $-$ and scaling of durations; $\dfn{totalMonths}$, $\dfn{totalDays}$, $\dfn{totalSeconds}$ \\
Predicates & $\dfn{startsWith}$, $\dfn{endsWith}$, $\dfn{contains}$, $\dfn{matches}$ (a regular expression) \\
\bottomrule
\end{tabularx}
\end{center}

Adding months keeps the day of the month where it exists and takes the last
day of the month otherwise ($31$~January $+$ \code{P1M} is the last day of
February). Timezoned dates and times are computed in UTC; the difference of
an instant with a timezone and one without is undefined.

\begin{rationale}
\textbf{A fixed library, not model-defined functions.} NORMA has no built-in
functions: a model defines its own, untyped. \FBM\ instead fixes one typed
library over the sorts, so a rule can be type-checked, evaluated and written
as SQL. The NORMA importer reads a model's functions by symbol or name into
this library and refuses the ones it does not know, with a reason in the
import log; the exporter defines the ones a rule uses in the exported model.
\end{rationale}

% ═════════════════════════════════════════════════════════════════════════════
\section{Constraints as formulas}
\label{sec:constraints}

This section lists the formula of every kind of constraint, as
\code{formalise} builds it. $R$ is a fact type with roles $1..n$; for a set
$U$ of roles, $\bar x_U$ are the variables of those roles and $\bar y$ the
others. $Q_i$ is the query of the $i$-th role sequence of an external
constraint (Section~\ref{sec:queries}); $Q_i(\bar x)$ says that $\bar x$ is
an answer of it, and $Q_i(\bar x, t)$ that it is an answer for target $t$.
Readings are omitted here; Section~\ref{sec:notation} shows how the app
writes the same formulas.

\subsection{Internal constraints of a fact type}

\begin{center}
\small
\begin{tabularx}{\textwidth}{l X}
\toprule
\textbf{Constraint} & \textbf{Formula} \\
\midrule
uniqueness over $U \subsetneq$ roles & $\forall \bar x_U\, \exists^{\le 1} \bar y\, R(\bar x)$ \\
uniqueness over all roles & no formula: it holds for any set of facts \\
mandatory role $i$ & $\forall x_i\, \exists \bar y\, R(\bar x)$ \\
frequency over $U$, ranges $[\mathit{lo}_j,\mathit{hi}_j]$ & $\forall \bar x_U\, \bigl(\exists\bar y\, R(\bar x) \to \bigvee_j \exists^{\mathit{lo}_j..\mathit{hi}_j}\bar y\, R(\bar x)\bigr)$ \\
value range on role $i$ & $\forall x_i\,\bigl(\exists\bar y\, R(\bar x) \to \mathit{in}(\val(x_i))\bigr)$ \\
cardinality range on role $i$ & $\bigvee_j \exists^{\mathit{lo}_j..\mathit{hi}_j} x_i\, \exists \bar y\, R(\bar x)$ \\
\bottomrule
\end{tabularx}
\end{center}

In a frequency range a lower bound of 1 or less is implied by
$\exists\bar y\,R$, so it leaves $\exists^{\le \mathit{hi}}$; an open upper bound gives
$\exists^{\ge\mathit{lo}}$, and a single number $\exists^{=n}$. The same
simplifications apply to cardinality ranges, with 0 as the implied lower
bound; a range that every population meets has no formula. A value range is
a disjunction of conditions, one per part: $v \approx c$, $v \ge
\mathit{lo}$ ($>$ for an open bound), $v \le \mathit{hi}$ ($<$), or both;
for an entity type it is about the reference value, as $\val$ says, and a
type without a data value gets no formula. Order conditions are not allowed
on unordered sorts.

\paragraph{Preferred identifiers.} A preferred internal uniqueness
constraint over all roles but the entity's role $e$ also says that each
entity has exactly one identifying combination:
$\forall x_e\,\exists^{=1}\bar y\,R(\bar x)$. When an explicit uniqueness
constraint on $e$ or a mandatory role on $e$ gives one half, only the other
half is added (``at most one'' or ``at least one'').

\subsection{Object types}

\begin{center}
\small
\begin{tabularx}{\textwidth}{l X}
\toprule
\textbf{Constraint} & \textbf{Formula} \\
\midrule
subtype link $S\sqsubseteq T$ & $\forall x{:}S\; T(x)$ \\
value range of a value type $V$ (or entity type) & $\forall x{:}V\; \mathit{in}(\val(x))$ \\
cardinality range of a type $T$ & $\bigvee_j \exists^{\mathit{lo}_j..\mathit{hi}_j} x{:}T\; \top$ \\
\bottomrule
\end{tabularx}
\end{center}

\subsection{External constraints}

\begin{center}
\small
\begin{tabularx}{\textwidth}{l X}
\toprule
\textbf{Constraint} & \textbf{Formula} \\
\midrule
inclusive-or & $\forall t\, \bigvee_i Q_i(t)$ \\
exclusive-or & $\forall t\, \bigl(\bigvee_i Q_i(t) \wedge \bigwedge_{i<j}\neg(Q_i(t)\wedge Q_j(t))\bigr)$ \\
subset & $\forall \bar x\,\bigl(Q_1(\bar x) \to Q_2(\bar x)\bigr)$ \\
equality & $\forall \bar x\,\bigwedge_{j>1}\bigl(Q_1(\bar x) \leftrightarrow Q_j(\bar x)\bigr)$ \\
exclusion & $\forall \bar x\,\bigwedge_{i<j}\neg\bigl(Q_i(\bar x) \wedge Q_j(\bar x)\bigr)$ \\
external uniqueness & $\forall \bar x\,\exists^{\le 1} t\, Q(\bar x, t)$ \\
\quad preferred identifier & also $\forall t\,\exists^{=1}\bar x\, Q(\bar x, t)$ \\
external frequency & $\forall \bar x\,\bigl(\exists t\,Q(\bar x,t) \to \bigvee_j \exists^{\mathit{lo}_j..\mathit{hi}_j} t\, Q(\bar x, t)\bigr)$ \\
value comparison $\mathit{op}$ & $\forall t, x_1, x_2\,\bigl(Q(x_1,x_2,t) \to \val(x_1)\ \mathit{op}\ \val(x_2)\bigr)$ \\
\bottomrule
\end{tabularx}
\end{center}

All sequences of a set-comparison constraint must have the same length. A
value comparison's stored $=$ and $\ne$ become $\approx$ and $\not\approx$;
both sides need data values of the same sort, and an order operator needs an
ordered sort.

\paragraph{Outer joins.} An external uniqueness constraint may be read with
outer joins (Halpin's ``o'' through the bar): an object that lacks an
optional part of the combination counts as having ``no value'' there, and
two such objects must differ elsewhere. The positions fall into groups that
share only the target; for every set $A$ of optional groups there is one
formula, $\forall \bar x\,\exists^{\le 1} t\,\bigl(\bigwedge_{G\notin
A}P_G(\bar x,t) \wedge \bigwedge_{G\in A}\neg\exists\bar y\,P_G(\bar
y,t)\bigr)$, where $P_G$ is the group's part of the query and $A=\emptyset$ is
the ordinary formula. At least one group must be mandatory.

\subsection{Ring constraints}

For a ring constraint over $R(x,y)$ (a query of arity 2 over one type):

\begin{center}
\small
\begin{tabular}{l l}
\toprule
\textbf{Ring property} & \textbf{Formula} \\
\midrule
irreflexive & $\forall x\,\neg R(x,x)$ \\
locally reflexive & $\forall x,y\,(R(x,y)\to R(x,x))$ \\
purely reflexive & $\forall x,y\,(R(x,y)\to x=y)$ \\
asymmetric & $\forall x,y\,(R(x,y)\to\neg R(y,x))$ \\
symmetric & $\forall x,y\,(R(x,y)\to R(y,x))$ \\
antisymmetric & $\forall x,y\,(R(x,y)\wedge x\ne y\to\neg R(y,x))$ \\
transitive & $\forall x,y,z\,(R(x,y)\wedge R(y,z)\to R(x,z))$ \\
intransitive & $\forall x,y,z\,(R(x,y)\wedge R(y,z)\to \neg R(x,z))$ \\
acyclic & $\forall x\,\neg\tcl{v}{w}{R(v,w)}{x,x}$ \\
strongly intransitive & $\forall x,y,z\,\bigl(R(x,y)\wedge\tcl{v}{w}{R(v,w)}{y,z}\to\neg R(x,z)\bigr)$ \\
\bottomrule
\end{tabular}
\end{center}

\subsection{Deontic constraints}

Every alethic formula $\varphi$ above has a deontic counterpart
$\Ob(\varphi)$, for an external or formula constraint, an internal frequency, a
mandatory role, a (non-preferred) internal uniqueness constraint, a value
range and a cardinality range. A preferred identifier is always alethic.

\subsection{Objectification}

An objectified fact type $f(A_1,\dots,A_n)$ with link fact types $f_i(D,
A_i)$, $D$ the objectified type, has three formulas:
\begin{align*}
&\forall \bar a\,\bigl(f(\bar a) \to \exists^{=1}d\,\textstyle\bigwedge_i f_i(d,a_i)\bigr) && \text{each fact has exactly one nesting instance,}\\
&\forall \bar a\,\bigl(\exists d\,\textstyle\bigwedge_i f_i(d,a_i) \to f(\bar a)\bigr) && \text{linked values form a fact,}\\
&\forall d\,\exists^{=1}\bar a\,\textstyle\bigwedge_i f_i(d,a_i) && \text{each nesting instance links exactly one combination.}
\end{align*}

\subsection{Implicit constraints}

Two kinds of constraint have no shape on the diagram; the Constraints tab
lists them on request, and validation, accounting and the equivalence check
use them.
\begin{itemize}
  \item \textbf{Disjoint roots.} The roots of two different entity-type
        hierarchies have no instance in common: $\forall x{:}A\;\neg B(x)$.
        Value types are left out: their disjointness does not depend on the
        schema.
  \item \textbf{Implicit participation} (Halpin). An instance of an entity
        type $T$ that is not independent plays some role other than those of
        its preferred identifier: $\forall x{:}T\,\bigvee_k
        \exists\bar y\,R_k(\dots x\dots)$, over every role played by a type of
        $T$'s hierarchy; when the role's player $P$ is neither $T$ nor one of
        its supertypes, the disjunct is $\exists\bar y\,(P(x)\wedge R_k(\dots
        x\dots))$. The entity types of a relational schema count as
        independent, since a row may have only its key. A fully derived type
        whose rule has its subject play such a role needs no entry.
\end{itemize}

\subsection{Formula constraints}

A constraint of kind \emph{formula} is a closed formula of \FBM\ written by
the user in the rule editor, with an optional name and modality; its
formula is exactly what was written. It is the general form of which the
graphical constraints are shorthands, and FreeFBM recognises when a formula
constraint has the shape of a graphical one (a subset, equality, exclusion,
frequency or uniqueness, inclusive- or exclusive-or, or value comparison,
with join paths) and offers to turn it into that constraint.

\begin{rationale}
\textbf{Counting quantifiers before aggregation.} Uniqueness and frequency
could be written with $\dfn{Count}$: $\forall x\,(\dfn{Count}\{y:R(x,y)\}\le
1)$. \FBM\ uses $\exists^{\le n}$ instead, so every graphical constraint is a
first-order formula with counting quantifiers, which the verbalisation reads
as ``at most one'', the negative rewrite turns into ``more than one'', and
the chase and the SQL writer handle without aggregates. Aggregation is
reserved for rules that compute values.
\end{rationale}

% ═════════════════════════════════════════════════════════════════════════════
\section{Queries: join paths of external constraints}
\label{sec:queries}

Each role sequence of an external constraint denotes a \emph{query}, a
formula whose free variables are the sequence positions and, for
uniqueness, frequency, inclusive-or and exclusive-or, the target:
\[
  \{\ \mathit{free}: [q_1,\dots,q_k],\ \mathit{target}: t,\ \mathit{body}: \varphi\ \}.
\]
A sequence of roles of one fact type has the implicit query $R(\dots)$; a
single subtype link has the query $S(t)$. Any other sequence needs a join
path, which is a query like any other (\code{src/logic/query.js}). Queries
have the same format as derivation rules (Section~\ref{sec:derivation}),
and the same evaluator answers both.

\paragraph{Anchoring.} Every position $q_i$ must play the role of its
sequence member in an atom of the query's \emph{positive conjunctive core},
the atoms reached from the top through $\exists$ and $\wedge$ only (not under
$\neg$, $\vee$, $\to$, $\leftrightarrow$, $\forall$ or a counting quantifier);
a role of an objectified fact type may be played through its link fact
type. The target must occur in the query.

\paragraph{The graphical fragment.} A query whose body is $\exists$ and
$\wedge$ over fact atoms, type atoms of subtypes and object equalities, with
every position anchored, can be drawn and edited as a graph of atoms and
links; the formula is derived from the graph. Any other query is edited as a
formula. Both are stored with their formula, so a file always carries it.

\begin{rationale}
\textbf{One format for queries and rules.} Join paths used to be a separate
graph language. They are now formulas in the rule format, with the graph as
an editor for the part that is intuitively graphical. Anchoring is stricter
than NORMA, which allows constant and calculated projections: in \FBM\ such
a projection is a derived fact type, explicit and visible on the diagram.
\end{rationale}

% ═════════════════════════════════════════════════════════════════════════════
\section{Derivation rules}
\label{sec:derivation}

A derivation rule defines the population of a type or fact type by a
formula. A \emph{derived} element's population is exactly what its rule
gives (marked $*$); a \emph{semi-derived} element's population is its stored
population plus what its rule gives (marked $+$).

\subsection{Subtypes}

A rule for subtype $S$ is a formula $\varphi(x)$ with one free variable,
the \emph{subject}, over the intersection of $S$'s direct supertypes:
\[
  \text{derived: } \forall x\,\bigl(S(x) \leftrightarrow \varphi(x)\bigr)
  \qquad
  \text{semi-derived: } \forall x\,\bigl(\varphi(x) \to S(x)\bigr)
\]
The candidates are the instances all direct supertypes share, so a derived
subtype keeps its supertypes' identifiers. A rule that does not mention the
subject gets a warning: it holds for all candidates or for none.

\subsection{Entity types through their reference scheme}

An entity type $E$ without supertypes but with a reference mode or a
composite identifier can be derived too. Its rule's free variables are its
reference values $a_1,\dots,a_k$, one per part of the scheme, and each answer
is an instance, the one those values identify:
\[
  \text{derived: } \forall \bar a\,\bigl(\exists e\,\mathit{ident}(e,\bar a) \leftrightarrow \varphi(\bar a)\bigr)
  \qquad
  \text{semi-derived: } \forall \bar a\,\bigl(\varphi(\bar a) \to \exists e\,\mathit{ident}(e,\bar a)\bigr)
\]
where $\mathit{ident}$ is the conjunction of the identifying fact atoms. The
rule must mention every reference value, and the reference scheme may not
use a derived fact type.

\subsection{Fact types}
\label{sec:factrules}

A fact type $p$ has one variable per role. A role is either
\emph{selected}, its variable bound to an object of the population, or
\emph{computed}, its object being the one whose data value a term gives.
With $\bar x$ the selected variables, $\bar y$ the computed ones and $\bar z$
the rule's extra variables, a rule is
\[
  p(\bar x, \bar y^{*}) \;\Leftarrow\; \varphi(\bar x, \bar z) \;\wedge\; \textstyle\bigwedge_i y_i \approx t_i(\bar x, \bar z)
\]
and is formalised as
\begin{align*}
  \mathrm{F1}\quad & \forall \bar x, \bar z\,\bigl(\varphi \to \exists \bar y\,(p(\bar x,\bar y) \wedge \textstyle\bigwedge_i y_i\approx t_i)\bigr) && \text{every derived fact holds,}\\
  \mathrm{F2}\quad & \forall \bar x, \bar y\,\bigl(p(\bar x,\bar y) \to \exists \bar z\,(\varphi \wedge \textstyle\bigwedge_i y_i\approx t_i)\bigr) && \text{only derived facts hold.}
\end{align*}
A derived fact type has both, a semi-derived one only F1. Without computed
roles the two are shown together as
$\forall\bar x\,(p(\bar x)\leftrightarrow\exists\bar z\,\varphi)$ (semi-derived:
$\forall\bar x,\bar z\,(\varphi\to p(\bar x))$). A condition $\top$ is left
out, and a derived fact type without extra variables shows F1 as
$\forall\bar x\,(\varphi\to\exists\bar y\,p(\bar x,\bar y))$, since F2 then
fixes $y_i\approx t_i(\bar x)$ for every fact.

The conditions on such a rule:
\begin{itemize}
  \item $\varphi$ may not use the computed variables, and a term may not use
        a computed variable;
  \item a computed role's player is a value type or an entity type with a
        (possibly inherited) reference mode, neither derived nor
        objectified; the term's sort must fit its data value;
  \item a selected or extra variable used neither in $\varphi$ nor in a term
        gets a warning, since it then ranges over its whole type;
  \item objectified fact types and fact types that identify an entity type
        cannot be derived.
\end{itemize}
Evaluating the rule adds the objects it computes to the populations of their
types (and their supertypes, and the value type of their reference mode).

\begin{rationale}
\textbf{Why F1 and not the biconditional.} The natural translation of
``markup = retail price $-$ wholesale price'' is
$\forall\vr{Article},\vr{Money}\,(\mathit{hasMarkup}(\vr{Article},\vr{Money})
\leftrightarrow \dots \wedge \vr{Money}\approx m_1 \dop{-} m_2)$. But
$\forall\vr{Money}$ ranges over the \code{Money} objects that already exist,
so if no object has the value $4$, the formula says nothing about a markup
of 4. Computed roles make the creation of the object explicit, and F1 puts
the extra variables outside, so that several bindings of $\bar z$ with
different values for one $\bar x$ each give a fact rather than one
arbitrary one.
\end{rationale}

\subsection{Recursion}

A rule depends on the fact types and types its atoms use, on the types its
variables range over (unless a positive fact atom of the formula binds the
variable through a role of its own type), on the supertypes of a derived
subtype, and on the derived fact types that compute objects of a type it
depends on (\code{src/logic/derivationDependencies.js}). The dependency
graph must be acyclic, with one exception: a \emph{semi-derived} fact type
or subtype may use itself where it occurs \emph{positively}, under an even
number of negations, outside the left of $\to$, outside $\leftrightarrow$,
and under no counting quantifier other than $\exists^{\ge n}$. Its population
is then the least fixpoint of its rule over the stored population.

\begin{rationale}
\textbf{Transitive closure rather than general recursion.} Reachability is
what most schemas need recursion for (acyclic rings, ``reports to,
directly or indirectly''). $\TC$ gives it inside a single formula, with a
fixed meaning and a uniform reading (``by one or more steps''), without
stratification. The NORMA importer turns positive linear self-recursion in a
fully derived fact type into $\TC^{*}$, and the exporter turns it back. The
remaining general case, positive self-reference in semi-derived rules, has
the least-fixpoint meaning above. A closure is over binary relations of
single variables; the tuple-indexed form
$[\TC_{\bar v,\bar w}\,\varphi](\bar s,\bar t)$ of the draft is not
implemented.
\end{rationale}

% ═════════════════════════════════════════════════════════════════════════════
\section{Notation in the app}
\label{sec:notation}

The Constraints tab, the rule editor and every report write formulas the
same way (\code{src/logic/renderFormula.js}).
\begin{itemize}
  \item \textbf{Variables} are named after their type, with spaces removed
        and later words capitalised ($\vr{StaffMember}$). A subscript is added
        only when two variables of the same type can be in scope at once
        somewhere in the formula; otherwise the bare name is used. A
        variable over an intersection is named after its types:
        $\vr{Student\&Employee}$.
  \item \textbf{Fact atoms} are written in reading form, in parentheses,
        with the primary reading's words around the variables:
        $\fa{\vr{Person}\ \rd{lives in}\ \vr{City}}$. Type atoms are
        $\ta{Student}{\vr{Person}}$.
  \item \textbf{Quantifiers} are written one per variable, $\forall
        \vr{Person}\ \exists \vr{City}$; a counting quantifier carries its bound
        as a superscript and may bind several variables:
        $\exists^{\le 1}\vr{Course}, \vr{Term}$.
  \item \textbf{Constants} are quoted unless their sort is a number or a
        truth value: $\vr{Name} \approx \dc{'Smith'}$, $\vr{Credits} \ge \dc{5}$.
  \item \textbf{Copying} gives plain text with \verb|^{…}| for superscripts
        and \verb|_{…}| for the variables of $\TC$:
        \verb|[TC_{Course_2,Course_3} (Course_2 requires Course_3)](Course_1, Course_1)|
        (with subscript digits where this shows \verb|_2|).
\end{itemize}

For example, a uniqueness constraint over the first two roles of
\emph{Student enrols in Course in Term}, a mandatory role and an acyclic
ring:
\begin{align*}
  & \forall\vr{Student}\ \forall\vr{Course}\ \exists^{\le1}\vr{Term}\ \fa{\vr{Student}\ \rd{enrols in}\ \vr{Course}\ \rd{in}\ \vr{Term}} \\
  & \forall\vr{Person}\ \exists\vr{PersonName}\ \fa{\vr{Person}\ \rd{has}\ \vr{PersonName}} \\
  & \forall\vr{Course}_1\ \neg[\TC_{\vr{Course}_2,\vr{Course}_3}\ \fa{\vr{Course}_2\ \rd{requires}\ \vr{Course}_3}](\vr{Course}_1,\vr{Course}_1)
\end{align*}

\begin{rationale}
\textbf{Readings, not predicate names.} A formula in reading form can be
read by a modeller who knows the diagram. Arguments still follow the fact
type's role order, because the primary reading's slots follow it, so the
reading form and a predicate form $\mathit{enrolsIn}(x,y)$ are the same atom
written two ways.
\end{rationale}

% ═════════════════════════════════════════════════════════════════════════════
\section{Two representations}
\label{sec:representations}

The code holds formulas in two forms.

\begin{itemize}
  \item \textbf{Formula trees} (\code{src/logic/formula.js}) are what
        \code{formalise} produces and the Constraints tab, the verbalisation
        and the SQL writer read. Variables are objects compared by identity,
        with a type (or several); quantifiers and counts may bind several
        variables. They are built, never stored.
  \item \textbf{Rule trees} (\code{src/logic/derivationRule.js}) are what the
        model stores and the rule editor edits: derivation rules, formula
        constraints, queries, mapping patterns and table rules. Variables have
        ids and are declared one per quantifier; a part not filled in yet is a
        hole; an entity variable may be used directly as a data term. A
        fact type's rule is \code{\{ body, extraVars, terms \}}, a query
        \code{\{ free, target, body \}}.
\end{itemize}

\code{ruleFormula} turns a checked rule tree into a formula tree (adding the
reference-mode atoms of Section~\ref{sec:values}); \code{formulaToRule}
(\code{src/logic/formulaToRule.js}) turns a formula tree back into a rule
tree, so that a formalised constraint can be evaluated, rewritten or kept as
a formula constraint. A counting quantifier of several variables has no rule
form except for the bounds 0 and 1, where at least one is $\exists\bar y\,\varphi$
and at most one is $\neg\exists\bar y,\bar y'\,(\varphi\wedge\varphi'\wedge
\bar y\ne\bar y')$.

% ═════════════════════════════════════════════════════════════════════════════
\section{What FreeFBM does with formulas}
\label{sec:uses}

\paragraph{Verbalisation.} Every entry's formula is the only input of the
verbalisation pipeline. The negative mode states the counterexample: $\varphi$
is rewritten as $\neg X$ with the negation pushed inward ($\forall$ becomes
$\exists$, $\exists^{\le n}$ becomes $\exists^{>n}$, value comparisons flip)
and $X$ is said as ``it is impossible that $X$''; a deontic $\Ob(\varphi)$
becomes $\Fb(X)$. The default mode states missing mandatory and uniqueness
constraints as $\Diamond X$. See the verbalisation specification.

\paragraph{Population validation and derived populations.} One evaluator,
in \code{src/logic/derivedPopulation.js}, computes the populations of derived
elements and evaluates formula constraints, semi-derivation rules and
queries on the sample population, with counterexamples. A script
(\code{scripts/population-coverage-check.mjs}) checks that the built-in
population checks report a constraint exactly when its formula is false, on
the sample and on mutated populations.

\paragraph{Constraint accounting, equivalence and Fold.} The mapping tools
build the smallest violation of a formula, carry it across a mapping and
close it under the other schema's constraints by a chase, which applies
formulas of the form $\forall\bar y\,(L\to R)$ with positive $L$ and $R$ as
tuple-generating dependencies.

\paragraph{SQL.} A constraint of a relational schema without a table form is
written from its formula: a \code{CHECK} when it reads one row, else a query that
returns the violations, enforced by triggers or assertions per SQL system
(\code{src/relational/formulaSql.js}). SQL conditions are three-valued; the
translation keeps them two-valued in both directions: an imported \code{CHECK}
becomes $\forall t\,N(\varphi)$, with $N$ (``not false'') and $U$ (``not
true'') expressing SQL's treatment of NULL in \FBM\
(\code{src/relational/checkFormula.js}), and an exported condition is never
unknown.

\paragraph{NORMA.} Derivation rules and join paths are read from and written
to NORMA role paths where they are conjunctive (with negation, disjunction,
conditions and calculated values as NORMA allows), and kept as informal
rules holding the formula otherwise.

\subsection{Mappings}
\label{sec:mappings}

A mapping rule relates a pattern over one schema to a pattern over another.
Patterns are the fragment of \FBM\ with $\exists$, $\wedge$, fact and type
atoms, $\neg$, value equalities $x\approx\$y$ and $x\approx\dc{'c'}$ with
\emph{frontier variables} $\$y$ shared by both sides, table-row atoms
(shorthand for the column fact atoms of a relational schema) and optional
parts. Their meaning, key-safety and the accounting of constraints across a
mapping are defined in \code{docs/logic/schema-mappings.tex}. View mappings
are derivation rules of a target schema's elements whose formulas read a
source schema.

% ═════════════════════════════════════════════════════════════════════════════
\section{Limits and open points}
\label{sec:open}

\begin{itemize}
  \item $\TC$ is over pairs of single variables; tuple-indexed closures are
        not implemented. The Oracle writer refuses a closure with
        parameters, since Oracle cannot correlate a recursive \code{WITH}
        inside a subquery.
  \item A rule tree's counting quantifier binds one variable; counts of
        several variables come only from \code{formalise} (frequency and
        external constraints) and have a rule form only for the bounds 0
        and 1.
  \item Entity types with a composite identifier have no data value, so
        value comparisons and terms reach them only through their parts.
        Converting between a reference mode and a composite identifier
        (value functions) is undecided.
  \item Objectified fact types cannot be derived.
  \item There are no local definitions (``$\psi$ \textbf{where} $p :=
        \varphi$''); an auxiliary predicate is a derived fact type.
  \item $\Ob$ is a marker on a whole constraint, not a modal operator that can
        be nested; there is no deontic logic beyond it.
\end{itemize}

% ═════════════════════════════════════════════════════════════════════════════
\appendix
\section{Where the logic lives in the code}
\label{app:code}

\begin{center}
\small
\begin{tabularx}{\textwidth}{l X}
\toprule
\textbf{Module} & \textbf{Contents} \\
\midrule
\code{src/logic/formula.js} & formula trees and their constructors \\
\code{src/logic/renderFormula.js} & notation: tokens, variable names, copy form \\
\code{src/logic/constraintFormulas.js} & \code{formalise}: the entries of Section~\ref{sec:constraints}, implicit constraints, objectification, defaults \\
\code{src/logic/values.js} & data values of variables (\code{valueAccess}, \code{withValues}) \\
\code{src/logic/dataOperators.js} & sorts, operators, predicates, aggregates, typing, term evaluation \\
\code{src/utils/datatypeValues.js} & value kinds: parsing, equality, order \\
\code{src/logic/derivationRule.js} & rule trees, typing of formulas, subtype and entity-type rules \\
\code{src/logic/factDerivation.js} & fact type rules, F1 and F2 \\
\code{src/logic/derivationDependencies.js} & dependencies and cycles between rules \\
\code{src/logic/query.js} & join-path queries, anchoring, the graphical fragment, outer joins \\
\code{src/logic/formulaConstraint.js} & formula constraints \\
\code{src/logic/formulaToRule.js} & formula trees back to rule trees \\
\code{src/logic/modality.js} & where deontic modality is stored \\
\code{src/logic/derivedPopulation.js} & the evaluator \\
\code{src/utils/independence.js} & independent types, identified roles \\
\code{src/verbalisation/negative.js} & counterexample rewrite \\
\bottomrule
\end{tabularx}
\end{center}

\end{document}
